\subsection{What is settled}

The discussion yields the following durable conclusions.
\begin{enumerate}[label=\arabic*.,leftmargin=*]
  \item The over-regularized handoff and the one-sided $\omega=1$ local tail
  are complete in the note-scoped KKT convention; the tail costs
  $O(1/(\rho\sqrt\alpha))$.
  \item A matched activation-tax proximal subproblem is individually local and
  can be solved locally at a constant inner contraction rate.
  \item The same linear activation tax prevents uniform exponent-two Bregman
  acceleration across changing support.  The corresponding dual envelope is
  not globally convex.
  \item A Euclidean proximal subproblem centered at a lower certificate is
  sandwiched below the RPPR optimum and has support volume at most $1/\rho$.
  \item For classical inertial Euclidean centers, cumulative active volume is
  controlled exactly by~\eqref{eq:volume-via-flux}.  Support growth itself is
  not the obstruction; repeated mass deactivation is.
  \item Branching trees support outward activation/deactivation waves.  The
  tree mechanism rules out using an unguarded $O(\alpha T)$ flux assumption in
  a graph-uniform proof.
  \item The product lower bound
  $\Omega(1/(\rho\sqrt\alpha))$ is rigorous for persistent-support one-hop
  methods, but not for arbitrary moving-frontier local algorithms.
\end{enumerate}

Accordingly, the unconditional graph-uniform work bound
\begin{equation}
  \Work_{\mathrm{hybrid}}
  =\softO\!\left(\frac1{\rho\sqrt\alpha}\right)
  \label{eq:still-open-target}
\end{equation}
must remain open for the classical extrapolated warmup.  It should not be
promoted into the active manuscript.

\subsection{The most credible safeguarded algorithm}

The tree calculation suggests a concrete redesign rather than another global
potential.  Maintain two states:
\begin{itemize}[leftmargin=*]
  \item an accelerated trial state, allowed to extrapolate; and
  \item a monotone lower shadow $\ell_t$, updated only by certified positive
  residual coordinate steps.
\end{itemize}
During an acceleration epoch, monitor a local wave certificate such as
\begin{equation}
 \mathfrak D_{a:b}
 :=\pos{\beta_{\mathrm A} p_{b-1}-p_b}
   +\beta_{\mathrm A}\sum_{t=a}^b
     p_{A_{t-2}\setminus A_t}(x_{t-2}),
 \label{eq:epoch-flux}
\end{equation}
or the weighted negative part suggested by~\eqref{eq:transformed-wave}.
When the certificate exceeds a fixed fraction of the source mass injected in
the epoch, discard momentum, retract to the lower shadow, enlarge the certified
face if necessary, and restart.

The observable theorem to seek is not that restarts never occur.  It is an
amortized statement of the form
\begin{equation}
  \Work_{\mathrm{total}}
  =\softO\!\left(
      \frac1{\rho\sqrt\alpha}
      +\frac{N_{\mathrm{rst}}}{\rho}
    \right),
  \label{eq:restart-target}
\end{equation}
where $N_{\mathrm{rst}}$ is the number of certified wave restarts, followed by
either a structural bound on $N_{\mathrm{rst}}$ or an instance-adaptive
guarantee that exposes it in the theorem. This formulation is compatible with
\cref{thm:bounded-core,thm:controlled-faces} and with classical adaptive
restart ideas \citep{odonoghue2015adaptive}, while using a graph-local signal
instead of a global objective test.

\subsection{Central missing lemmas}

\begin{description}[style=nextline,leftmargin=3.4cm]
  \item[Finite-tree excitation.] Complete the rigorous transfer from the
  outward period-11 wave to the zero-start root-source response on a sequence
  of finite simple trees.  This would convert the current computational and
  limiting evidence into a formal asymptotic counterexample to
  \cref{conj:flux}.

  \item[Restart amortization.] Prove that a wave-triggered restart pays for
  itself through objective decrease, newly certified KKT slack, or expansion
  of a bounded core.  Counting restarts without such a payment simply
  reintroduces the unknown $K$ in~\eqref{eq:face-phase-work}.

  \item[Lower-shadow accuracy.] Bound the gap between the accelerated trial
  point and the monotone lower shadow tightly enough that restarting does not
  lose the $1/\sqrt\alpha$ progress already obtained.

  \item[Flattened bell.] Determine whether
  $\softO((1/\sqrt\alpha)\log(1/\rho))$ guarded epochs can enforce
  $O(1/\rho)$ active volume per epoch.  The tree wave shows that a simple
  fixed momentum schedule cannot be the entire proof.

  \item[Broader oracle lower bound.] Either construct a hard family that
  forces simultaneous volume on a moving-frontier evolving-set algorithm, or
  prove that such an algorithm can evade the product lower bound.  Output size
  and iteration complexity alone do not settle this question.

  \item[Residual reconciliation.] Translate the note-scoped KKT threshold,
  the AESP degree-normalized infinity error, the APPR threshold, and the
  proximal fixed-point residual into one accepted repository convention before
  promoting any theorem.
\end{description}

\subsection{Final research position}

The hybrid idea remains viable, but the proof target has changed.  The
correct algorithm cannot be ``classical acceleration plus locality for free.''
It must make locality an enforced invariant or an observable restart
condition.  The strongest current theorem is the flux-dependent
bound~\eqref{eq:hybrid-flux-bound}, together with the bounded-core and
controlled-face corollaries.  The strongest obstruction is the combination
of the exact face-change defect~\eqref{eq:face-change-defect} and the rooted
tree wave~\eqref{eq:radial-obstacle}.  Any next proof should state explicitly
which of these two mechanisms it suppresses.
